\section{La Microsoft Kinect}

\subsection{Description}
La Kinect, connue à l'origine sous le nom "Project Natal" est une
périphérique informatique créé par Microsoft, comportant une caméra
couleur, un étalage de 4 microphones, des moteurs pour s'orienter avec 
accéléromètres 
et un système senseur de profondeur composé d'un 
projecteur infra-rouge et d'un capteur. 
  \begin{figure}[h!]
  \centering
  \includegraphics[width=0.8\linewidth]{images/kinect_diagram}
  \caption{Diagramme montrant les sous-systèmes de l'appareil Kinect.}
  \end{figure}
L'appareil n'utilise pas, comme on pourrait croire, le "Temps de Vol" pour
calculer la profondeur d'une scène. C'est en fait la déformation de chacune des 
ondes infra-rouge envoyée, lorsqu'elles frappent une surface, qui est mesuré. Cette technique dite de 
"Lumière Structurée"
lui permet de générer une image 3D d'une scène à une résolution de 640x480 avec 
une profondeur de 11 bits, qui sera mise à jour 30 fois par seconde.

\subsection{Historique}
Sorti en Novembre 2010, le système était prévu à l'origine pour la console
de jeu de Microsoft, la "Xbox360". Comme la plupart des consoles et
périphériques de jeu elle
est vendue à perte, l'idée étant de récupérer l'argent perdu en vendant
des jeux. Ceci en fait l'appareil de son genre le moins cher et le plus couramment
disponible.

Il est donc peu surprenant que, dès sa sortie, des pilotes Libres et gratuits
aient été créés afin de permettre le détournement du périphérique par des
chercheurs et des amateurs. Microsoft a répondu en début de 2011 avec un
premier SDK pour Windows initialement peu complet. Des mises à jours et une
deuxième version, sortie en Octobre 2012, comblent ces lacunes\cite{wiki_kinect}.

\subsection{Pilotes et bindings bas-niveau}

Il est important de noter que la reconnaissance vocale et l'extraction de 
squelette ne sont pas faits par l'appareil, ou "serveur", lui-même~: ce sont des
bibliothèques sur la console ou l'ordinateur, le "client", qui s'en chargent. 
Il faut donc bien
différencier les pilotes des bibliothèques de bas niveau, qui permettent 
d'accéder aux données brutes, des
bibliothèques de haut-niveau, qui interprètent ces données.

Le jour de la sortie de la Kinect, le 4 Novembre 2010, Adafruit Industries 
offre une prime de \$2000 (augmentée ensuite jusqu'à \$3000) pour le premier qui
sortira des pilotes permettant l'utilisation de la Kinect sur ordinateur\cite{adafruit_bounty}.

\subsubsection{Libfreenect}
\begin{figure}[h!]
\centering
\includegraphics[width=\linewidth]{images/hector}
\caption{Première sortie de la Libfreenect par son créateur, Hector.}
\end{figure}
Le prix est remporté 6 jours plus tard le 10 Novembre par un programmeur Linux
connu sous le nom "Hector"\cite{adafruit_winner}. C'est la première version de Libfreenect autour
duquel le groupe OpenKinect se forme. Libfreenect a pour avantage son
très grand nombre de bindings~: C/C++, Java, Python, Actionscript, Javascript,
LISP, C\#, pour ne nommer qu'un sous-ensemble. 
Il est toujours maintenu (dernier commit Décembre 2012).
\begin{figure}[h!]
\centering
\includegraphics[width=0.5\linewidth]{images/openkinect_logo}
\caption{Logo de la OpenKinect.}
\end{figure}

\subsubsection{CL NUI}
9 jours après la sortie de Libfreenect la compagnie Coding Laboratories 
(CL) et le groupe logiciel-libre
Natural User-Interaction (NUI) sortent CL-NUI\cite{clnui}. Nous sommes alors 2 semaines 
après la sortie de l'appareil. CL-NUI propose des bindings 
WPF\footnote{Windows Presentation Foundation}/C\# C et C++ et permet d'accéder 
à tout sauf le microphone. Le développement continue jusqu'en fin Décembre 2010 
mais,
le défi déjà remporté, s'arrête par la suite. Notons pourtant qu'il faudra 
attendre 2 ans pour avoir accès aux données des accéléromètres à travers la 
SDK officiel\ldots
\subsubsection{Kinect for Windows}  
Les pilotes officiels de Microsoft viennent avec leur SDK, "Kinect for 
Windows". Inutile de dire
qu'ils ne fonctionnent que sous Windows. Ce qui est un peu plus surprenant, 
c'est qu'il n'y a de soutient que pour Window 7 et 8, donc pas pour Windows XP 
ou Vista.
\subsubsection{SDK OpenNI}
OpenNI (Open Natural Interaction) est un consortium à but non lucratif composé
des grands acteurs économiques du domaine. S'y figure notamment PrimeSense, la 
compagnie qui a développé la Kinect pour Microsoft au premier lieu. Leur SDK
est plus ou moins analogique à celui de Microsoft où il est question de
Kinect, mais fonctionne aussi sur les plus anciennes 
versions de Windows mais
aussi sous Linux (y compris Android) et Mac OSX.
\begin{figure}[h!]
\begin{minipage}{0.49\linewidth}
  \centering
  \includegraphics[width=0.9\linewidth]{images/openni_logo}
\end{minipage}
\begin{minipage}{0.49\linewidth}
  \centering
  \includegraphics[width=0.9\linewidth]{images/primesense_logo}
\end{minipage}

\caption{Logos de la OpenNI et de PrimeSense.}
\end{figure}
La OpenNI permet également de s'abstraire du matériel utilisé. Il suffit que
les pilotes utilisés soient conformes à ses standards~: pour passer, par 
exemple, de la Kinect à une caméra PrimeSense nous devons simplement remplacer 
les pilotes "SensorKinect" avec "PrimeSensor".
\begin{figure}[h!]
\centering
\includegraphics[width=0.6\linewidth]{images/primesense_camera}
\caption{Caméra Carmine 1.09 de PrimeSense.}
\end{figure}

% ------------------------------------------------------------------------------  

\subsection{Middleware}
Les bibliothèques dites "middleware interaction naturelle" servent ici en 
générale à faire une 
extraction de squelette et une reconnaissance de gestes. Elles permettent donc 
au développeur de s'abstraire des données brutes. Le groupe OpenNI définit 
des normes auxquelles doit obéir cette couche, de 
manière à faciliter la communication à travers la pile de technologies (voir 
figure~\ref{fig:openni_technology_stack}).
\begin{figure}[h!]
\centering
\includegraphics[width=0.7\linewidth]{images/openni_technology_stack}
\caption{Pile de technologie d'après la OpenNI.}
\label{fig:openni_technology_stack}
\end{figure}


\subsubsection{Kinect for Windows}  
Bien que les SDK de Microsoft et de la OpenNI signifient tout deux 
"Simple Development Kit",
celui de Microsoft, "Kinect for Windows", contient à la fois les pilotes, 
l'API de bas niveau et une couche d'analyse des données qu'OpenNI appellerait 
middleware. Il n'est donc pas fait pour interfacer avec les couches de la 
OpenNI et ne suit par leurs normes.

Kinect for Windows a quelques avantages par rapport à NITE~: il propose un
placement prédictif des articulations et génère deux points en plus :
le poignet et le bout des doigts. Cela étant, sa prédiction consomme beaucoup 
plus de ressources et peut amener à de faux positifs\cite{microsoft_vs_openni}.

\subsubsection{NITE de PrimeSense}
NITE (Natural Interface Technology for End-User) est un framework Libre et 
multiplateforme développé par PrimeSense. Vu que c'est PrimeSense qui a 
lancé OpenNI, NITE est devenu la tête d'affiche de ses
middleware standardisés. Elle se repose donc sur la SDK de la OpenNI. 
\begin{figure}[h!]
\centering
\includegraphics[width=0.3\linewidth]{images/nite_logo}
\caption{Logo du framework NITE.}
\end{figure}

\subsubsection{OpenCV}
OpenCV (Open Computer Vision) est une vaste bibliothèque de toutes sortes de 
fonctions aillant un rapport avec la détection de formes à partir d'images
statiques ou dynamiques. Elle existe depuis 1999, donc bien avant la Kinect.
\begin{figure}[h!]
\centering
\includegraphics[width=0.4\linewidth]{images/opencv_logo}
\caption{Logo de la OpenCV.}
\end{figure}
Étant donnée qu'il s'agit d'une bibliothèque généraliste, elle a peu de 
fonctionnalités qui ciblent précisément les données 3D, contrairement aux 
bibliothèques dites "d'interaction naturelle".

% ------------------------------------------------------------------------------  

\subsection{Bibliothèques de haut-niveau}

\subsubsection{Zigfu}
Crée en Mai 2011, Zigfu est une petite entreprise formée de 4 personnes dont 
Amir Hirsch, le fondateur, et 2 ex-employés de PrimeSense. Il fournit, 
par dessus de l'extraction de squelette et de la reconnaissance
de gestes, une API permettant de créer des boutons, listes, curseurs et autres
éléments GUI, tous contrôlés par le mouvement.
\begin{figure}[h!]
\centering
\includegraphics[width=0.3\linewidth]{images/zigfu_logo}
\caption{Logo de Zigfu.}
\end{figure}
Zigfu intègre également la cinématique inverse et un système de 
prédiction et
d'interpolation afin d'avoir un squelette répondant à des contraintes 
physiques à 60Hz\cite{zigfu_video}. La plupart des jeux modernes tourne à cette 
vitesse et, 
rappelons que la caméra ne fournit que des données 30Hz.

Zigfu fonctionne avec Kinect for Windows ou OpenNI avec NITE. Il est 
avant-tout orienté Web, avec des
bindings HTML 5 et Unity, et une version Flash en cours de développement
\cite{zigfu_review}.

\subsubsection{FAAST}
La Flexible~Action~and~Articulated~Skeleton~Toolkit (FAAST) est une 
bibliothèque qui se repose soit sur Kinect for Windows, soit OpenNI et NITE, 
comme Zigfu. Elle n'est pour l'instant compatible qu'avec Windows.

Développé par la Interactive~Media~Division (MxR) de 
l'Institute~of~Creative~Technologies (ICT) de 
l'University~of~Southern~California (USC). Elle permet de monter à un niveau 
encore plus abstrait~: l'utilisateur travaille avec des événement générés à
partir de postures et des gestes, au lieu d'analyser manuellement le 
squelette.

% ------------------------------------------------------------------------------  

\subsection{Résumé}
Le schéma \ref{fig:technology_overview} résume les interactions entre les diverses 
technologies que nous ayons vues.

\begin{figure}[h!]
\centering
\includegraphics[width=0.9\linewidth]{images/technology_overview}
\caption{Pile de technologie de la Kinect.}
\label{fig:technology_overview}
\end{figure}
